\emph{Notation.} $\Cl(1,3)$ with $\gamma_{0}^{2} = 1$, $\gamma_{j}^{2} = -1$ ($j = 1, 2, 3$), $I_{4} = \gamma_{0}\gamma_{1}\gamma_{2}\gamma_{3}$, $I_{3} = \gamma_{1}\gamma_{2}\gamma_{3} = \gamma_{0}I_{4}$, $\sigma_{j} = \gamma_{j}\gamma_{0}$, $I\sigma_{3} = I_{4}\sigma_{3} = \gamma_{2}\gamma_{1}$. Spinors are elements of $\Cl^{+}(1,3)$; $L_{a}\psi = a\psi$, $R_{b}\psi = \psi b$. $I_{4}$ commutes with every even element and anticommutes with every vector, so on spinors $L_{I_{4}} = R_{I_{4}}$, and $R_{b}$ commutes with every $L_{a}$. Two basis blades either commute or anticommute.

\emph{Step 1 (the null kernel; (i)).} Let $k = \omega(\gamma_{0} + n)$, $\omega > 0$, $n \cdot \gamma_{0} = 0$, $Q(n) = -1$, and $B_{k} = \gamma_{0}n$. Then $k = \omega\gamma_{0}(1 + B_{k})$ and $B_{k}^{2} = -\gamma_{0}^{2}n^{2} = 1$, so $k\psi = 0$ if and only if $B_{k}\psi = -\psi$: $S_{k}$ is the $-1$ eigenspace of the involution $L_{B_{k}}$ on $\Cl^{+}(1,3)$. For a unit transverse $t$ ($t \cdot \gamma_{0} = t \cdot n = 0$, $Q(t) = -1$) the even element $u = \gamma_{0}t$ anticommutes with $B_{k}$, so $L_{u}$ exchanges the two eigenspaces and $\dim S_{k} = 4$. With $D_{k} = I_{4}B_{k}$ and $\psi \in S_{k}$, $D_{k}\psi = I_{4}B_{k}\psi = -I_{4}\psi$, so $L_{D_{k}} = -L_{I_{4}}$ on $S_{k}$. For $k = \gamma_{0} + \gamma_{1}$, $B_{k} = \gamma_{0}\gamma_{1}$ and $D_{k} = \gamma_{2}\gamma_{3}$: the wedge boost and dual plane of \claimref{synthesis.wedge-reflections-compose-to-the-two-pi-rotor-so-geometric-modular-conjugations-give-interacting-spin-statistics}, whose identity $I_{4}B = D$ is the one-line computation above (re-derived here, not taken as a premise). The argument holds for every $n$ at once; equivalently, a rotor $R$ fixing $\gamma_{0}$ carries $\gamma_{0} + \gamma_{1}$ to every future null direction, and $(RkR^{-1})(R\psi) = Rk\psi$ shows $RS_{k} = S_{RkR^{-1}}$, so the case $k = \gamma_{0} + \gamma_{1}$ suffices.

The plane $D_{k}$ is orthogonal to $\gamma_{0}$ and $n$, and $R_{\theta} = \exp(\tfrac{\theta}{2}D_{k})$ is the rotation by $\theta$ about $n$, right-handed for the spatial orientation $I_{3} = \gamma_{0}I_{4}$ (for $n = \gamma_{1}$: $R_{\theta}\gamma_{2}R_{\theta}^{-1} = \gamma_{2}e^{-\theta\gamma_{2}\gamma_{3}} = \cos\theta\,\gamma_{2} + \sin\theta\,\gamma_{3}$). On $S_{k}$, $R_{\theta}\psi = (\cos\tfrac{\theta}{2} - \sin\tfrac{\theta}{2}\,I_{4})\psi = e^{-I_{4}\theta/2}\psi$: writing the action of a one-parameter group as $e^{-I_{4}\theta J}$, the spin about $n$ is $J = +\tfrac12$ on all of $S_{k}$. For the mode $\psi(x) = e^{-sI_{4}k\cdot x}\psi_{0}$, $\psi_{0} \in S_{k}$, the translation by $a$ gives $\psi(x - a) = e^{sI_{4}k\cdot a}\psi(x)$, so with $e^{I_{4}a\cdot P}$ the energy-momentum relative to $I_{4}$ is $P = sk$: energy $s\omega$, momentum direction $sn$. The rotation about $sn$ by $\theta$ is the rotation about $n$ by $s\theta$, with generator $sJ$: the helicity relative to $I_{4}$ is $s/2$. This restates, for use, \claimref{synthesis.light-ray-good-part-of-a-spinor-is-two-chirality-arrays}(iii)-(iv) (there with the boost bivector $n\gamma_{0} = -B_{k}$ of the generator, which fixes the same $S_{k}$) and \claimref{field-theory.where-each-lives}(4), and agrees with clause (iii) of \claimref{field-theory.fock-space-maxwell-modes}.

\emph{Step 2 (the chirality; (ii)).} $\Gamma = L_{I_{4}}R_{I\sigma_{3}}$ has $\Gamma^{2} = L_{I_{4}^{2}}R_{(I\sigma_{3})^{2}} = 1$ and $\Gamma R_{I\sigma_{3}} = L_{I_{4}}R_{(I\sigma_{3})^{2}} = -L_{I_{4}}$, which is (ii). $\Gamma$ commutes with $L_{a}$ for even $a$ and anticommutes with $L_{v}$ for vectors $v$; so it commutes with the left rotor action, and $k\Gamma\psi = -\Gamma k\psi$ shows that it preserves $S_{k}$. $R_{\sigma_{1}}$ commutes with $L_{k}$ and anticommutes with $\Gamma$ ($\sigma_{1}$ anticommutes with $I\sigma_{3} = \sigma_{1}\sigma_{2}$), so it exchanges the two eigenspaces of $\Gamma$ on $S_{k}$, each of dimension $2$. The projectors $\chi_{e} = \tfrac12(1 - e\Gamma)$ are those of the check of \claimref{field-theory.fock-space-maxwell-modes} (built there with $I\sigma_{3} = \gamma_{2}\gamma_{1}$), so its label $e$ is $-\Gamma$. On a mode of chirality $\Gamma = \gamma$, (ii) gives $R_{I\sigma_{3}} = -\gamma L_{I_{4}}$: the charge rotation $\psi \mapsto \psi e^{I\sigma_{3}\alpha}$ acts as $e^{-\gamma I_{4}\alpha}$, and $e^{-sI_{4}k\cdot x}\psi_{0} = \psi_{0}e^{-s_{H}I\sigma_{3}k\cdot x}$ with $s_{H} = -\gamma s$: the frequency sign relative to $I_{4}$ is minus the chirality times the frequency sign relative to the right unit $R_{I\sigma_{3}}$. This is \claimref{field-theory.where-each-lives}(4) in this notation.

\emph{Step 3 (covariant maps; (iii)).} For a bivector blade $B$, $1 + tB$ is a positive multiple of a rotor ($B^{2} = -1$, any real $t$; $B^{2} = 1$, $|t| < 1$), so the linear span of the rotors contains $1$ and every bivector; being closed under products it is $\Cl^{+}(1,3)$, which the bivectors generate. A real-linear $T$ commuting with every left rotor multiplication therefore commutes with every $L_{a}$, $a \in \Cl^{+}(1,3)$, and $T(\psi) = T(L_{\psi}1) = L_{\psi}T(1) = R_{T(1)}\psi$: $T$ is right multiplication by $T(1)$, and it commutes with $L_{I_{4}}$. (In the language of the centre: $\Cl^{+}(1,3)$ is a simple algebra with centre $\mathrm{span}\{1, I_{4}\}$, and the commutant of its left regular action is its right action; the elementary form above needs no structure theory.) Such a map preserves every $S_{k}$ and every frequency sign $s$, since $T(e^{-sI_{4}k\cdot x}\psi_{0}) = e^{-sI_{4}k\cdot x}T(\psi_{0})$.

\emph{Step 4 (the discrete maps; (iv)).} The four maps $\Lambda$ act on vectors by $x$, $\gamma_{0}x\gamma_{0}$, $-\gamma_{0}x\gamma_{0}$, $-x$; each is orthogonal and its own inverse, so $k\cdot\Lambda^{-1}x = (\Lambda k)\cdot x$. For the mode $\psi$ above, $a\psi(\Lambda^{-1}x)b = a\,e^{-sI_{4}(\Lambda k)\cdot x}\psi_{0}b = e^{-s(-1)^{|a|}I_{4}(\Lambda k)\cdot x}\,a\psi_{0}b$, since an odd blade anticommutes with $I_{4}$. If $\Lambda k$ is past, write $k' = -\Lambda k$. So the image is the mode with future $k' = \epsilon\Lambda k$, $\epsilon = \sgn(\gamma_{0}\cdot\Lambda k)$, frequency sign $s' = s(-1)^{|a|}\epsilon$ and amplitude $a\psi_{0}b$; $\epsilon = +1$ for $\Lambda = 1, P$ and $-1$ for $T, PT$.

\emph{Which maps keep the equation.} A mode solves the massless equation exactly when its amplitude lies in the null kernel: $\nabla(e^{-sI_{4}k\cdot x}\psi_{0}) = -sk I_{4}e^{-sI_{4}k\cdot x}\psi_{0} = sI_{4}e^{sI_{4}k\cdot x}k\psi_{0}$. So the map keeps the equation on all modes if and only if $(\Lambda k)a\psi_{0} = 0$ for every future null $k$ and every $\psi_{0} \in S_{k}$ ($b$ is invertible). For $\Lambda = 1, PT$ put $c = a$; for $\Lambda = P, T$, $(\Lambda k)a = \pm\gamma_{0}k(\gamma_{0}a)$, put $c = \gamma_{0}a$; the condition is $kc\psi_{0} = 0$. Write $ckc^{-1} = \tilde k$, the vector $k$ with the signs of its components along the generators that anticommute with $c$ reversed; the condition is $\tilde k\psi_{0} = 0$ on $S_{k}$. If $c \in \{1, I_{4}\}$ (up to sign), $\tilde k = \pm k$ and it holds. Otherwise the grade $r$ of $c$ is $1$, $2$ or $3$, so some generator lies in $c$ and some does not, and these two classes are reversed with opposite signs ($\gamma_{\mu}$ in $c$ is multiplied by $(-1)^{r-1}$, $\gamma_{\mu}$ not in $c$ by $(-1)^{r}$); one class contains $\gamma_{0}$ and the other some $\gamma_{j}$, and for $k = \gamma_{0} + \gamma_{j}$, $\tilde k = \pm(\gamma_{0} - \gamma_{j})$, whose kernel is the $+1$ eigenspace of $L_{\gamma_{0}\gamma_{j}}$, while $S_{k}$ is its $-1$ eigenspace (Step 1): $\tilde k\psi_{0} \ne 0$ for $\psi_{0} \ne 0$. So the family is exactly $a \in \pm\{1, I_{4}\}$ for $\Lambda = 1, PT$ and $a \in \pm\{\gamma_{0}, \gamma_{0}I_{4}\}$ ($\gamma_{0}I_{4} = I_{3}$) for $\Lambda = P, T$, with $b$ any blade of the parity of $a$: $16$ maps up to sign for each $\Lambda$. Hence $(-1)^{|a|}\epsilon = +1, -1, +1, -1$ for $\Lambda = 1, P, T, PT$: the frequency sign flips exactly for $P$ and $PT$, whatever $a$ and $b$, that is, exactly for the maps that reverse the spatial orientation $I_{3}$ (both $P$ and $T$ reverse $I_{4}$; $T$ also reverses $\gamma_{0}$, and $I_{3} = \gamma_{0}I_{4}$). The reversal under which Premise \claimref{synthesis.n-plus-null-ground-state} calls the filled sea the time reverse of the null vacuum is the $I_{4}$-linear one, $PT$ here ($a$ even), which flips $s$; the $T$ of this family has odd $a$, is $I_{4}$-antilinear, and keeps $s$.

\emph{Chirality, helicity and charge of the image.} $\Gamma(a\psi_{0}b) = I_{4}a\psi_{0}bI\sigma_{3} = (-1)^{|a|}c(b)\,a(\Gamma\psi_{0})b$, where $c(b) = \pm 1$ as $b$ commutes or anticommutes with $I\sigma_{3}$; $\Gamma$ commutes with the phase $e^{-s'I_{4}k'\cdot x}$, so the image of a chirality-$\gamma$ mode has chirality $\gamma' = (-1)^{|a|}c(b)\gamma$. By Step 1 its helicity is $s'/2$, and by Step 2 its charge sign relative to $I_{4}$ is $-\gamma'$. Examples: $a = 1$, $b = \sigma_{1}$, $\Lambda = 1$ (the spacetime-algebra charge conjugation $\psi \mapsto \psi\sigma_{1}$): $s' = s$ and $c(\sigma_{1}) = -1$, so the chirality flips and the frequency sign does not; $a = b = \gamma_{0}$, $\Lambda = P$ (parity): both flip; $a = \gamma_{0}$, $b = \gamma_{0}\sigma_{1} = -\gamma_{1}$, $\Lambda = P$ (their product, a CP map): $s' = -s$ and $\gamma' = (-1)(-1)\gamma = \gamma$, so the frequency sign flips and the chirality does not. With the identity these realise all four patterns, and by Step 2 every pair $(s, \gamma)$ occurs on every $S_{k}$: the two labels are independent.

\emph{Step 5 (the consequence; argued).} The datum of Premise \claimref{synthesis.n-plus-null-ground-state}, which of the null vacuum and the filled sea is the state, is fixed relative to the unit $I_{4}$ and the time orientation, and on massless modes it is the choice of which frequency sign $s$ relative to $I_{4}$ counts as positive energy. By Step 4 this sign is odd under every map of the family that reverses the spatial orientation and under CP, and even under the charge conjugation; the chirality $\Gamma$ is odd under the charge conjugation and parity and even under CP. A rule identifying the datum with $\Gamma$ would equate a bit odd under CP with a bit even under it, so it is invariant under none of these maps; and by (iii) no pointwise Lorentz-covariant map changes $s$ at all. Since the helicity relative to $I_{4}$ is $s/2$ for either chirality (Step 1), the datum can equally be read as a helicity orientation of the massless modes; it is not the chirality $\Gamma$. That $\Gamma$ is the handedness of the series' doublets, the sign of $I_{6}$ relative to $I_{4}$ on $L_{0}$, $L_{\infty}$ (\claimref{field-theory.where-each-lives}(3)), is assumed, not shown; whether some other structure of the series ties the datum to the doublets is not addressed. The checks verify Steps 1 to 4 exactly on sampled null directions (rotation and boost images of $\gamma_{0} + \gamma_{1}$), the commutant dimension, and the full list of blade pairs.
